\documentclass[11pt]{article}
\usepackage{amsmath,amssymb,amsthm}
\usepackage[margin=1in]{geometry}

\newtheorem{theorem}{Theorem}
\newtheorem{lemma}[theorem]{Lemma}

\title{Problem 516: 5-Smooth Totients}
\author{Project Euler Solutions}
\date{}

\begin{document}
\maketitle

\section{Problem Statement}

A number is \emph{5-smooth} (Hamming number) if its largest prime factor does not exceed~5. Let $S(L)$ be the sum of all $n \leq L$ with $\varphi(n)$ being 5-smooth. Given $S(100) = 3728$, find $S(10^{12}) \bmod 2^{32}$.

\section{Mathematical Foundation}

\begin{theorem}[Characterization]
A positive integer $n$ has $\varphi(n)$ being 5-smooth if and only if
\[
n = 2^a \cdot 3^b \cdot 5^c \cdot p_1 p_2 \cdots p_k,
\]
where $a, b, c \geq 0$, $k \geq 0$, and each $p_i > 5$ is a distinct prime with $p_i - 1$ being a Hamming number.
\end{theorem}

\begin{proof}
$(\Leftarrow)$\; By multiplicativity of $\varphi$:
\[
\varphi(n) = \varphi(2^a)\,\varphi(3^b)\,\varphi(5^c)\prod_{i=1}^{k}(p_i - 1).
\]
Each factor is 5-smooth by hypothesis, and a product of 5-smooth numbers is 5-smooth.

$(\Rightarrow)$\; If $q > 5$ is prime and $q^2 \mid n$, then $q \mid \varphi(q^2) = q(q-1)$ divides $\varphi(n)$, so $\varphi(n)$ has a prime factor $q > 5$, contradicting 5-smoothness. Thus each prime $q > 5$ dividing $n$ appears to the first power, and $\varphi(q) = q - 1$ must itself be 5-smooth.
\end{proof}

\begin{lemma}[Hamming Number Count]
The number of Hamming numbers $\leq L$ is $O(\log^3 L)$. For $L = 10^{12}$, there are approximately 3{,}428.
\end{lemma}

\begin{proof}
A Hamming number $2^a 3^b 5^c \leq L$ requires $a \leq \log_2 L$, $b \leq \log_3 L$, $c \leq \log_5 L$. The triple count is $O(\log^3 L)$.
\end{proof}

\begin{lemma}[Hamming Prime Count]
The number of primes $p \leq L$ with $p - 1$ being 5-smooth is at most $O(\log^3 L)$. For $L = 10^{12}$, there are approximately 545.
\end{lemma}

\begin{proof}
Each such prime corresponds to a distinct Hamming number $p - 1$, so the count is bounded by the Hamming number count.
\end{proof}

\section{Algorithm}

\begin{enumerate}
    \item Generate all Hamming numbers $h \leq L$ by enumerating $2^a \cdot 3^b \cdot 5^c$.
    \item Identify Hamming primes: test $h + 1$ for primality for each Hamming number $h$.
    \item Enumerate valid $n$ via DFS: at each level, either include the next Hamming prime (at most once) or stop. Multiply the partial product of Hamming primes by each Hamming number $\leq L/\mathrm{product}$ and accumulate.
    \item Return the sum modulo $2^{32}$.
\end{enumerate}

\section{Complexity Analysis}

\begin{itemize}
    \item \textbf{Hamming generation:} $O(\log^3 L)$.
    \item \textbf{Primality tests:} $O(\log^3 L \cdot \sqrt{L})$ worst case (fast with Miller--Rabin).
    \item \textbf{DFS enumeration:} Pruned by $\leq L$ bound; runs in under 1 second for $L = 10^{12}$.
    \item \textbf{Space:} $O(\log^3 L)$ for the Hamming list.
\end{itemize}

\section{Answer}

\[
\boxed{939087315}
\]

\end{document}
